% frontmatter/polish/mapa-del-curso.tex -- el año entero en una tabla
% (frontmatter/mapa-del-curso.tex, en polaco; el plan completo, semana a
% semana, está en notes/01-plan.md).
\section*{Plan roku}

\begin{center}
\renewcommand{\arraystretch}{1.4}
\begin{tabularx}{\linewidth}{@{}c c c c X@{}}
\toprule
\textbf{Część} & \textbf{Dni} & \textbf{Pora roku} & \textbf{Pieniądze} & \textbf{Czego się uczymy} \\
\midrule
1 & 1--65 & jesień & do 20 € & Wartość rzeczy, handel wymienny i
pieniądze; kiedy nie starcza dla wszystkich; to, co wspólne; potrzeby i
zachcianki; oszczędzanie; cena. \\
2 & 66--130 & zima & do 50 € & Wybieranie; dawne pieniądze; kto
produkuje, a kto kupuje; dbanie o rzeczy; budżet przyjęcia; zrobić
samemu czy kupić; usługi; oszczędzanie na cel. \\
3 & 131--195 & wiosna & do 100 € & Dawne ceny; ile kosztuje przepis;
naprawić czy kupić nowe; zawody i pensja; sklepy w okolicy; skarbonka;
od kwiatka do słoika; bank; kiermasz. \\
4 & 196--260 & lato & do 100 € & Najpierw to, co najważniejsze;
porównywanie cen; mierzenie i ważenie; zapisywanie rachunków; cena
miodu; kieszonkowe na festyn; to, co ma wartość, choć nie kosztuje
pieniędzy; stoisko z lemoniadą; planowanie. \\
\bottomrule
\end{tabularx}
\end{center}

\vspace{6mm}
Każda część ma 13 tygodni po 5 dni -- także w czasie wakacji: zeszyt to
nadal jedna strona dziennie, od poniedziałku do piątku, w Boże
Narodzenie i latem. Każdy tydzień ma jedną myśl i jedno słowo, a
pieniądze rosną razem z rokiem: do 20 € jesienią, do 50 € zimą i do
100 € wiosną i latem, zawsze w całych euro.
\newpage
